% ----------------------------------------------------------------------
%  Iteration 3 --- drop-in section for research-paper-synthesis-graph.tex
%  Requires only the existing preamble (enumitem, hyperref, xcolor).
%  Tables are plain `tabular`; no booktabs needed.
% ----------------------------------------------------------------------
\section{Iteration 3}

Iteration 2 asked whether a typed knowledge graph retrieves better than
vector search. It answered \emph{no}, and diagnosed its worst case ---
relational queries at 0.202 --- as missing \texttt{undercuts} edges.
Iteration 3 changes the question from \emph{what does the system know} to
\emph{what does the system do}: decompose a query into sub-questions,
retrieve for each part, critique the evidence against the plan, re-retrieve
the gaps, and synthesise. It then asks whether acting beats the static arms
on the same gold set.

\subsection{Design/Coding Process}

\subsubsection*{The claim was fixed before the measurement}

The acceptance criterion and the \emph{suspicious} outcome were both written
down before the arm was run. A uniform gain across all query types was
pre-registered as evidence \emph{against} the hypothesis, on the reasoning
that decomposition issues more retrievals and more retrieval helps
everything --- so a gain that appeared everywhere would indicate volume
rather than planning. Two statistical tests were pre-registered, which fixes
the multiple-comparison correction in advance rather than choosing it once
the p-values are visible.

This matters because the result landed exactly on that pre-registered shape:

\begin{center}
\begin{tabular}{lrrr}
\hline
 & $\Delta$ vs static & W / L / T & $p$ \\
\hline
all 34 queries          & $+0.059$ & 12 / 9 / 13 & 0.403 \\
relational ($n=14$)     & $+0.250$ & 8 / 1 / 5   & \textbf{0.012} \\
non-relational ($n=20$) & $-0.075$ & 4 / 8 / 8   & 0.340 \\
\hline
\end{tabular}
\end{center}

The overall figure is a null. The effect lives entirely in the relational
subset, and the arm is measurably \emph{worse} everywhere else --- which is
what makes it credible rather than disappointing. Cost was $1.2\times$
(\$0.028 against \$0.024 per query), not the order of magnitude the plan had
anticipated.

\subsubsection*{Three properties carry the comparison's validity}

The arm sits behind the existing \texttt{System} protocol, so the evaluation
runner scored it without modification. Three construction decisions exist
purely so that the comparison isolates one variable:

\begin{enumerate}[label=\arabic*.]
  \item \textbf{Synthesis is the static arm's code, not a copy of it.} The
        agentic arm holds a \texttt{VectorRAGSystem} instance and calls its
        evidence formatting, handle resolution and synthesis directly, so the
        prompt and the citation-handle scheme are identical \emph{by
        construction} rather than by inspection. A test pins the function
        identity, so a future edit cannot silently fork them. The variable
        under test is the control flow and nothing else.
  \item \textbf{The retrieval budget binds.} Six retrievals per query,
        \emph{refused} rather than logged when exceeded, and at
        \texttt{top\_k = 20} rather than 60 --- so several small retrievals
        stay within the same order of evidence as one large static retrieval.
        A budget that warns instead of refusing is not a budget.
  \item \textbf{It fails closed.} A planner that raises, or returns an empty
        plan, degrades to a single retrieval on the original query and sets a
        \texttt{planner\_failed} flag. Without this a broken planner scores as
        a working one, because the fallback path still produces an answer.
\end{enumerate}

\subsubsection*{The determinism problem, and why the fix had to be architectural}

Three sampling defects were found before anything was measured, each the same
defect in a different layer, and each had been turning a reported figure into
a single draw:

\begin{enumerate}[label=\arabic*.]
  \item The \textbf{judge} ran at the provider default temperature (1.0).
        Nothing in the project had ever set one. Judging the same 34 answers
        three times with an identical rubric produced $\kappa$ spreads up to
        0.25, against a 0.26 gap to the trust threshold --- so the noise was
        the same size as the decision.
  \item \textbf{Extraction} had the same defect, meaning the entire graph was
        unreproducible. Extracting the same 21 papers twice from one prompt
        differed on 9 of 147 field outcomes.
  \item The \textbf{planner} is nondeterministic and pinning does not fix it.
\end{enumerate}

The third case is the interesting one. \texttt{gpt-5.4-nano} does not return
identical text at temperature 0, and in an agentic arm \emph{the planner's
output is the retrieval query}. A reworded sub-question embeds differently,
reaches different chunks, and changes which papers can be cited. A static arm
embeds a fixed gold query and absorbs that noise; this arm amplifies it.
Measured: \texttt{must\_cite\_recall} of 0.567 and 0.333 across two runs of
functionally identical code.

No temperature setting can fix a component whose output is prose that must
then be embedded. The fix was therefore structural --- the arm now always
retrieves on the \textbf{original query} first, so the plan moves the margin
of the evidence set rather than the whole of it:

\begin{center}
\begin{tabular}{lccc}
\hline
 & run 1 & run 2 & run 3 \\
\hline
with anchor    & 0.417 & 0.417 & 0.417 \\
without anchor & 0.367 & 0.383 & 0.417 \\
\hline
\end{tabular}
\end{center}

Spread $0.050 \rightarrow 0.000$. The change was adopted on a stopping rule
stated \emph{before} it was made --- keep it only if it narrows the spread ---
so the decision could not be rationalised after seeing the numbers.

The corpus was then rebuilt (271 papers, extraction pinned at temperature 0).
That moved the substrate further than intended: \texttt{Claim} nodes $+44\%$,
nodes $23{,}689 \rightarrow 27{,}777$, edges $11{,}131 \rightarrow 14{,}978$.
Iteration 2's stored figures do not describe this substrate, which is stated
rather than quietly absorbed.

\subsubsection*{The required ablation}

Removing the self-critique drops relational performance from 0.528 to 0.437 ---
back among the static arms ($+0.091$, 8 wins / 1 loss / 5 ties, $p = 0.043$).
Under Bonferroni across the two pre-registered tests ($\alpha = 0.025$) this
does \emph{not} survive, while the relational-versus-baseline result
($p = 0.012$) does. Both are reported at their corrected significance rather
than at the uncorrected one.

The critique's contribution is also \emph{unexplained}: it changes the
evidence on 31--33 of 34 queries but adds a \emph{required} paper on only
4--7. A step that rarely reaches new required papers should not be worth
0.091, and the trajectory measures do not say why it is.

\subsubsection*{The proposed mechanism was refuted}

Iteration 2 proposed that decomposition works by \emph{substituting} for the
missing \texttt{undercuts} edges --- the second hop becomes a second query
instead of a graph traversal. Three independent measurements test that
account and none supports it:

\begin{enumerate}[label=(\alph*)]
  \item \textbf{Adding the edges does not work.} Rewriting the contradiction
        prompt with the twelve labelled spurious edges as worked negatives held
        edge precision at $32.5\%$ --- identical stratum rates under both
        prompts --- while discarding roughly $65\%$ of the real edges.
  \item \textbf{Edge coverage was not the binding constraint.} The rebuild took
        \texttt{undercuts} from 33 to 119 and traversal from \emph{zero}
        occurrences across 34 queries to six. The typed-graph arm scored
        $0.383 \rightarrow 0.383$, unchanged to three decimal places. The gap
        was real; closing it changed nothing.
  \item \textbf{The gain does not concentrate where the account predicts.}
        Splitting the relational queries on whether their gold claims name a
        limitation or cost gives $+0.167$ where one is named against $+0.296$
        where none is --- the gain is \emph{larger} where the account says it
        should be smaller. Pre-registered as underpowered at 5/9 before the
        numbers were seen, so it cannot refute the account, but it does not
        support it.
\end{enumerate}

Separately, \emph{plan quality does not predict outcome}. Across three
independent runs, decomposition specificity correlates with
\texttt{must\_cite\_recall} at $\rho = -0.16, -0.08, -0.16$
($p = 0.375, 0.654, 0.369$). How distinctive a plan is carries no information
about whether the query is answered well.

What survives as an untested hypothesis: the gain comes from issuing several
retrievals against the \emph{parts} of a question, not from the plan being
good. Volume alone cannot explain it, since the arm is worse on
non-relational queries while spending the same retrievals there. So it
appears to be the act of splitting rather than the craft of the split. An arm
that splits mechanically on conjunctions, with no planner at all, would
separate the two --- and is carried forward.

\subsubsection*{Query-time writes, and an invariant that was never enforced}

The agent persists its decomposition as \texttt{STAGED} nodes with
provenance: the query, its sub-questions, and per sub-question the papers it
reached. That last field records which papers \emph{jointly} bear on one part
of a question --- something per-paper extraction structurally cannot observe,
because it never holds two papers at once.

Building it surfaced a latent safety defect. The store interface had promised
since Iteration 1 that metrics query \texttt{CURATED} only. Neither retrieval
arm actually filtered on the layer. The promise held solely because nothing
had ever written a \texttt{STAGED} node --- it was \emph{vacuously} true, and
the first query-time write would have made it false, letting an agent raise
its own score by writing to the graph it is scored against. Acceptance was
demonstrated with real staged data rather than argued:

\begin{center}
\begin{tabular}{lrr}
\hline
 & filtered & unfiltered \\
\hline
before staging & 20{,}688 & 20{,}688 \\
after staging 38 nodes & 20{,}688 & 20{,}726 \\
\hline
\end{tabular}
\end{center}

Promotion from staged to curated is now gated behind an explicit approval
flag.

\subsubsection*{One deliverable kept with its justification refuted}

Local inference was planned on the assumption of a loop issuing 5--10 calls
per query and an affordability problem. The measured arm makes \textbf{three}
calls at \$0.028 per query, and the whole deliverable cost about \$6. There
was no cost problem to solve, so the stated rationale is refuted.

The item was kept anyway, because its acceptance criterion is a
\emph{portability} claim that nothing else in the iteration addresses. The
plumbing is landed --- vLLM serves the OpenAI API, so ``route calls to vLLM''
reduces to ``point an OpenAI client elsewhere''. What waits is the run, and
the blocker is hardware rather than software: the target 14B model needs
about 8.5\,GB in 4-bit with CUDA-only kernels, against an 8\,GB machine whose
usable working set caps near 5.3\,GB. Recording the refuted justification
separately from the surviving goal keeps the plan auditable.

\subsection{Terms}

\begin{enumerate}[label=\arabic*.]
  \item \textbf{Planner--critic loop} --- an agentic control flow in which the
        system first plans (decomposes a query into sub-questions), acts
        (retrieves per sub-question), then critiques its own evidence against
        the plan and re-retrieves what is missing, before synthesising. The
        contrast is a \emph{static} arm: one retrieval, one synthesis, no
        self-inspection.
  \item \textbf{Anchored retrieval} --- always retrieving on the original,
        unmodified query in addition to the planner's sub-questions. Introduced
        because a planner's output \emph{is} a retrieval query, so its
        nondeterminism otherwise propagates into which papers are reachable.
        The anchor confines plan variation to the margin of the evidence set.
  \item \textbf{Fail closed} --- degrading to a safe, clearly-marked fallback
        when a component fails, rather than producing output that is
        indistinguishable from success. Here: a failed planner falls back to a
        single retrieval and sets a flag, so a broken planner cannot be scored
        as a working one.
  \item \textbf{CURATED / STAGED layers} --- the provenance split on every
        graph node and edge. \texttt{CURATED} is offline, reviewed material and
        is the only layer metrics may read; \texttt{STAGED} is anything the
        agent writes at query time, and is never merged automatically.
        Promotion between them requires explicit approval.
  \item \textbf{Trajectory evaluation} --- scoring the \emph{process} rather
        than the answer: whether the plan was well-formed, how efficiently
        retrieval reached required papers, what the critique changed. Its
        purpose is to explain \emph{why} an arm wins, which an outcome metric
        cannot.
  \item \textbf{Decomposition specificity} --- whether a plan's sub-questions
        match their own query's facets better than they match any other query's
        facets. Replaced a naive coverage measure that read 1.000 on every
        facet; measuring the null showed 100\% of \emph{unrelated} query--plan
        pairs also passing, so the original measure was detecting nothing.
  \item \textbf{Retrieval efficiency} --- required papers reached per retrieval
        call. Reported with its definitional overlap stated: it counts papers
        \emph{reached} while the headline metric counts papers \emph{cited}, so
        the two share a numerator term and part of their correlation is built
        in.
  \item \textbf{Pre-registration} --- writing the acceptance criterion, the
        planned statistical tests, and the outcomes that would count as
        \emph{evidence against} the hypothesis, before running the experiment.
        It is what allows a null overall result with a positive subgroup result
        to be reported as a finding rather than as a subgroup fished out
        afterwards.
  \item \textbf{Bonferroni correction} --- dividing the significance threshold
        by the number of pre-registered tests, so that running several tests
        does not inflate the chance that one passes by luck. At two tests the
        threshold becomes $\alpha = 0.025$.
  \item \textbf{Quadratic-weighted $\kappa$} --- agreement between two raters
        on an ordinal scale, corrected for agreement expected by chance, and
        penalising disagreements by the \emph{square} of their distance so that
        near-misses cost little. Preferred over correlation because the criteria
        are 1--5 levels: a judge that ranks answers identically but sits a point
        low has high correlation and mediocre $\kappa$, and those two failures
        need different fixes (rescaling versus a rubric rewrite).
  \item \textbf{Test--retest reliability} --- having one grader re-grade the
        same answers after an interval, giving a ceiling on how well \emph{any}
        automated judge could agree with them. Used here because the project has
        a single grader, so a second-annotator measure is permanently
        unavailable; it established that one criterion was already at the human
        ceiling and no rubric could have improved it.
  \item \textbf{Temperature} --- the sampling parameter controlling randomness
        in a model's output. A provider default of 1.0 makes every reported
        figure one draw from a distribution; pinning it to 0 is a precondition
        for a number being a measurement rather than a sample.
  \item \textbf{vLLM} --- a high-throughput local inference server that exposes
        an OpenAI-compatible API, so switching a client from a hosted provider
        to local hardware is a base-URL change rather than a rewrite.
\end{enumerate}

\subsection{Fixes}

Twenty-eight defects were found and fixed during the iteration. One pattern is
worth naming ahead of the list: \textbf{seven of them produced a plausible
number rather than an error}, and in four cases the thing that caught a later
mistake was a fix that had looked like housekeeping when it was made. A defect
that crashes is cheap; a defect that returns 0.000, or 32.5\%, or ``all facets
covered'', is expensive precisely because nothing signals it.

\subsubsection*{Sampling was never disabled (7)}

\begin{enumerate}[label=A\arabic*.]
  \item The judge ran at the provider default temperature. Found by judging the
        same 34 answers three times with one rubric: $\kappa$ spread to 0.25
        against a 0.26 gap to the trust bar.
  \item Extraction had the same defect, so the graph itself was
        unreproducible. Two extractions of 21 papers differed on 9 of 147 field
        outcomes. Pinned \emph{before} the corpus rebuild, so the rebuild
        happened once.
  \item The planner is nondeterministic and pinning cannot fix it, because its
        output is the retrieval query. Fixed architecturally by the anchor.
  \item Criteria were being certified from a single draw. A \texttt{-{}-repeats}
        mode was added, and certification now requires the \emph{worst} sample
        to clear the bar --- certifying on the best or median draw is what makes
        a lucky sample look trustworthy.
  \item The judge's per-criterion justifications were requested and then
        discarded. This is the only reason a later rubric defect was
        diagnosable: the saved reasons showed it penalising a correctly reused
        citation handle. A score cannot explain itself.
  \item There was no way to re-score stored answers, so a rubric change could
        not be isolated from resampled synthesis. A re-judge tool now re-grades
        answers on disk and never modifies the source run.
  \item Rubric versions were not retained, leaving no control arm. Version 1 was
        kept and made selectable --- and then \emph{scored best of three}, so the
        default was set back to it, with a test recording why.
\end{enumerate}

\subsubsection*{Gold-set handling (3)}

\begin{enumerate}[label=B\arabic*.]
  \item Scripts defaulted to the 10-query file even for a 34-query run, which
        would have scored a third of the data and reported it as the whole. Now
        auto-selects the smallest covering file, and errors if none covers.
  \item Calibration figures depended on which gold set was used, unstated --- and
        the sets disagree on which criteria pass. Calibration now always prints
        the complement subset beside the one it calibrated on, so a figure can
        never be quoted without its denominator.
  \item The evaluation entry point hardcoded the active gold file. A
        \texttt{-{}-gold} flag was added and non-default runs are tagged in the
        run directory name.
\end{enumerate}

\subsubsection*{The reproducibility layer --- three faults, not one (3)}

\begin{enumerate}[label=C\arabic*.]
  \item The artifact node type was reachable from only two section types, while
        four of the five gold papers that state a code URL state it elsewhere.
  \item The field-list hint was gated on a different node type, so a section
        could be asked for the node type while being told nothing about which
        attributes to fill.
  \item The schema had no member for the phrasing 15 corpus papers actually
        use, which made two authored gold values \emph{unscoreable by any
        extraction} --- a ceiling of 4/6 before the extractor was involved.
\end{enumerate}

Result: two fields moved from zero correct to two each; total correct 8 to 11.
The item had been carried forward as a single routing gap; it was three
independent faults, only one of which was the named one.

\subsubsection*{Contradiction tooling (4)}

\begin{enumerate}[label=D\arabic*.]
  \item The verdict cache key recorded the claim pair but not the prompt
        version, so a version-2 run would have returned 16,965 \emph{version-1}
        verdicts and reported them as version 2's. One command away from a false
        finding.
  \item The scoring path read the accepted-edge total from the old edge file
        regardless of what it was auditing, reporting ``of 3,072 accepted'' for
        a pass that accepted 1,172.
  \item Input and output files were hardcoded, so auditing a second attempt
        would have destroyed the record of the first.
  \item The labeller used for the original audit had never itself been
        validated. A validation mode was added and the labeller \emph{immediately
        failed} it: 76.7\% agreement, but only 2 of 15 real disagreements found
        --- a labeller that answers ``neither'' to everything reports low
        precision regardless of the truth. Without this check the number would
        have been reported.
\end{enumerate}

\subsubsection*{A safety invariant that was never enforced (3)}

\begin{enumerate}[label=E\arabic*.]
  \item The typed-graph arm did not filter on layer --- in \emph{two} places,
        seeding and expansion, either of which leaves a route in.
  \item The citation-graph arm had the same gap in both places.
  \item Promotion from staged to curated had no approval gate, so unreviewed
        material could enter the scored layer. It now refuses without an
        explicit flag.
\end{enumerate}

Stated as an honest limit: the acceptance experiment exercised the \emph{node}
filter; the edge filter is a second lock and was not exercised, because staged
edges currently connect only to staged nodes.

\subsubsection*{Trajectory measures (2)}

\begin{enumerate}[label=F\arabic*.]
  \item Gold identifiers carry a \texttt{paper:} prefix while the trace records
        bare identifiers, so the intersection was empty for every query and
        retrieval efficiency read \textbf{0.000 as a plausible number rather
        than an error}. The strongest argument for building the trajectory
        measures before the experiment that consumes them.
  \item Decomposition coverage read 1.000 on all 29 facets. Measuring the null
        --- each facet against every \emph{other} query's plan --- gave 0.769
        against 0.808 matched, with 100\% of unrelated pairs passing at the
        default threshold. Replaced by a paired specificity measure plus a
        validity check that prints before the scores.
\end{enumerate}

\subsubsection*{Plumbing (5)}

\begin{enumerate}[label=G\arabic*.]
  \item Per-paper flags on the extraction entry point, which prices a routing
        change at about \$0.36 instead of \$6.72.
  \item An optional trajectory field on the system output plus runner
        persistence; static arms leave it empty so their traces stay
        byte-identical and remain directly comparable.
  \item The report writer made public, so the re-judge tool renders its means
        identically --- a re-judged run that reported differently would not be
        comparable to the run it re-scores.
  \item Base-URL and temperature threaded through the LLM layer and all three
        clients.
  \item Gold-set resolution moved into the schema module, since two scripts need
        it and the scripts directory is not importable.
\end{enumerate}

\subsubsection*{Documentation that contradicted the measurements (2)}

\begin{enumerate}[label=H\arabic*.]
  \item The README certified one judge criterion and rejected another --- both
        backwards --- and carried a stale query count, node count, test count and
        ``next iteration''. Its opening still asserted the typed-graph bet as
        untested after two iterations had tested it.
  \item The co-citation section explained how it had made the typed-versus-citation
        ablation fair, and never recorded that the ablation was subsequently
        lost.
\end{enumerate}
